\section*{Chapitre \ref{ch:b1:nucleophilic-substitution} --- La substitution nucléophile}

\begin{solution}{exo:b1:nucleophilic-substitution:1}
\ce{CH3CH2OH + Br-}~; \ce{CH3OCH2CH3 + I-}~; \ce{CH3CH2CH2CN + Cl-}~;
\ce{CH3NH3+ + Br-}. \omterm{def:b1:nucleophilic-substitution:substitution}{Substrats}~: les halogénoalcanes~; \omterm{def:b1:electronic-effects:nucleophile}{nucléophiles}~:
\ce{OH-}, \ce{CH3CH2O-}, \ce{CN-}, \ce{NH3}~; \omterm{def:b1:electronic-effects:nucleophile}{groupes partants}~: les ions
halogénure.
\end{solution}

\begin{solution}{exo:b1:nucleophilic-substitution:2}
Ordre 2~: vitesse triplée~; les deux divisées par deux, vitesse divisée
par 4. Ordre 1 par rapport à RY seulement~: vitesse inchangée~; divisée
par deux.
\end{solution}

\begin{solution}{exo:b1:nucleophilic-substitution:3}
SN2 (primaire, bon \omterm{def:b1:electronic-effects:nucleophile}{nucléophile}, aprotique)~; SN1 (tertiaire, eau)~; SN2
(méthyle, \omterm{def:b1:electronic-effects:nucleophile}{nucléophile} fort)~; \omterm{def:b1:nucleophilic-substitution:sn1}{solvolyse} SN1 (secondaire, \omterm{def:b1:electronic-effects:nucleophile}{nucléophile}
faible et neutre, \omterm{def:b1:intermolecular-forces-solvents:solvent-classes}{solvant protique}), éventuellement accompagnée d'un peu
de SN2.
\end{solution}

\begin{solution}{exo:b1:nucleophilic-substitution:4}
Le cyanure prend la place opposée au brome et les trois autres groupes
basculent. Dans le produit, le groupe CN est de rang 1, comme l'était
Br, les autres rangs étant inchangés~: la disposition inversée porte le
descripteur opposé, \cip{S}-2-méthylbutanenitrile.
\end{solution}

\begin{solution}{exo:b1:nucleophilic-substitution:5}
1. \omterm{def:b1:electronic-effects:cleavage}{Hétérolyse}~: flèche de la liaison C--Br vers Br, qui donne
\ce{(CH3)3C+} et \ce{Br-} (lente). 2. Un \omterm{def:b1:lewis-resonance-vsepr:lewis}{doublet non liant} de l'eau vers
le carbone cationique~: \ce{(CH3)3C-OH2+}. 3. Un \omterm{def:b1:lewis-resonance-vsepr:lewis}{doublet non liant} d'une
autre molécule d'eau capte un proton du groupe \ce{OH2+}, le doublet de
O--H revenant sur l'oxygène. Seule la première étape, lente, fixe la
vitesse~; l'eau, solvant, est en si grand excès que sa concentration ne
varie pas.
\end{solution}

\begin{solution}{exo:b1:nucleophilic-substitution:6}
Le carbone qui porte Br est primaire, mais il est voisin d'un groupe
\textit{tert}-butyle qui occupe l'espace situé derrière la liaison
C--Br~: l'\omterm{def:b1:elementary-steps:energy-profile}{état de transition} de la SN2 est très encombré. La SN1
exigerait un \omterm{def:b1:electronic-effects:intermediates}{carbocation} primaire, trop instable.
\end{solution}

\begin{solution}{exo:b1:nucleophilic-substitution:7}
L'iodure remplace l'iodure par SN2, chaque fois avec inversion~:
\cip{S} devient \cip{R} jusqu'à ce que les deux soient présents en
quantités égales, un \omterm{def:b1:stereochemistry-in-depth:enantiomers}{mélange racémique}. Chaque substitution transforme
une molécule \cip{S} en une molécule \cip{R}~: la rotation diminue de
l'équivalent de deux molécules, si bien qu'elle décroît deux fois plus
vite que le \omterm{def:b1:nucleophilic-substitution:substitution}{substrat} n'est substitué.
\end{solution}

\begin{solution}{exo:b1:nucleophilic-substitution:8}
SN2~: bromométhane $>$ 1-bromopropane $>$ 2-bromopropane $>$
2-bromo-2-méthylpropane (encombrement). SN1~: l'ordre inverse
(stabilité du \omterm{def:b1:electronic-effects:intermediates}{carbocation}), le bromométhane et le 1-bromopropane ne
réagissant presque pas.
\end{solution}

\begin{solution}{exo:b1:nucleophilic-substitution:9}
$x_{\text{inv}} - x_{\text{rét}} = 0.12$ et $x_{\text{inv}} + x_{\text{rét}} = 1$~:
\qty{56}{\percent} d'inversion, \qty{44}{\percent} de rétention.
Essentiellement une \omterm{def:b1:nucleophilic-substitution:sn1}{racémisation} SN1, avec un léger excès d'inversion~:
le bromure qui s'en va masque encore son côté quand l'eau arrive.
\end{solution}

\begin{solution}{exo:b1:nucleophilic-substitution:10}
$v = k_1k_2[\mathrm{RY}][\mathrm{Nu}]/(k_{-1}[\mathrm{Y^-}] + k_2[\mathrm{Nu}])$
(\cref{prop:b1:nucleophilic-substitution:sn1-law}). Ajouter \ce{Y-}
augmente le dénominateur~: davantage de \omterm{def:b1:electronic-effects:intermediates}{carbocations} reviennent au
\omterm{def:b1:nucleophilic-substitution:substitution}{substrat} au lieu d'être capturés, et la vitesse diminue. Effet
négligeable lorsque
$k_2[\mathrm{Nu}] \gg k_{-1}[\mathrm{Y^-}]$, comme lorsque le \omterm{def:b1:electronic-effects:nucleophile}{nucléophile}
est le solvant.
\end{solution}

\begin{solution}{exo:b1:nucleophilic-substitution:11}
Concentrations égales~: $-\mathrm{d}c/\mathrm{d}t = kc^2$, $1/c = 1/C_0 + kt$,
$t_{1/2} = 1/(kC_0) = \qty{5.0e5}{s}$ (environ six jours). Avec un excès
d'hydroxyde d'un facteur cinquante, $[\ce{OH-}] \approx 50C_0$ est
constante~: $c = C_0
e^{-k't}$ avec $k' = 50kC_0 = \qty{1.0e-4}{s^{-1}}$, $t_{1/2} = \ln 2/k' =
\qty{6.9e3}{s}$, environ deux heures.
\end{solution}

\begin{solution}{exo:b1:nucleophilic-substitution:12}
Le bromure devrait expulser \ce{OH-}, \omterm{def:b1:predominant-reaction:strong-acid}{base forte} et mauvais \omterm{def:b1:electronic-effects:nucleophile}{groupe
partant}~: pas de réaction. Dans HBr concentré, le groupe OH est protoné
en \ce{-OH2+}, dont le \omterm{def:b1:electronic-effects:nucleophile}{groupe partant} est l'eau. Alcool tertiaire~: l'eau
part, ce qui donne le \omterm{def:b1:electronic-effects:intermediates}{carbocation}, capturé par \ce{Br-} (SN1). Alcool
primaire~: \ce{Br-} attaque le carbone de \ce{R-OH2+} par l'arrière
pendant que l'eau s'en va (SN2).
\end{solution}

\begin{solution}{pb:b1:nucleophilic-substitution:1}
\textbf{1.} Chaque molécule qui réagit donne un \ce{H+} et un \ce{Cl-}~;
la conductivité est la somme des contributions ioniques, proportionnelle
à la quantité d'ions.
\textbf{2.} $[\mathrm{RCl}]/[\mathrm{RCl}]_0 = (\kappa_\infty -
\kappa)/\kappa_\infty$.
\textbf{3.} $\ln[\mathrm{RCl}] = \ln[\mathrm{RCl}]_0 - kt$~: tracer
$\ln(\kappa_\infty - \kappa)$ en fonction de $t$.
\textbf{4.} 0.693, 0.603, 0.513, 0.423, 0.333, 0.153, $-0.117$, $-0.387$.
\textbf{5.} Oui, de pente $-0.0180$ par minute partout~: ordre 1.
\textbf{6.} Un éventuel ordre par rapport à l'eau~: sa concentration est
constante (dégénérescence de l'ordre).
\textbf{7.} $k = 0.0180/60 = \qty{3.0e-4}{s^{-1}}$.
\textbf{8.} $t_{1/2} = \ln 2/k = \qty{2.3e3}{s}$, soit \qty{38.5}{min}.
\textbf{9.} $\ln 10/k = \qty{7.7e3}{s}$, soit \qty{128}{min}.
\textbf{10.} $\ln(\kappa_\infty - \kappa)$~: 0.693 à $t = 0$, 0.109 à
\qty{10}{min}, et ainsi de suite~: pente $-0.0584$ par minute,
$k = \qty{9.7e-4}{s^{-1}}$.
\textbf{11.} $9.7/3.0 = 3.2$.
\textbf{12.} À \qty{30}{min}~: $2.000(1 - e^{-0.0584 \times 30}) = 1.653$,
pour 1.654 mesuré.
\textbf{13.} SN1~: \omterm{def:b1:nucleophilic-substitution:substitution}{substrat} tertiaire, \omterm{def:b1:intermolecular-forces-solvents:solvent-classes}{solvant protique}, \omterm{def:b1:electronic-effects:nucleophile}{nucléophile}
neutre.
\textbf{14.} Comme dans l'\cref{exo:b1:nucleophilic-substitution:5}, avec
Cl.
\textbf{15.} L'ionisation lente ne fait intervenir que le \omterm{def:b1:nucleophilic-substitution:substitution}{substrat}~: $v =
k[\mathrm{RCl}]$.
\textbf{16.} L'hydroxyde ne pourrait intervenir que dans l'étape rapide
de capture, après l'ionisation cinétiquement déterminante~; la vitesse
n'en dépend pas.
\textbf{17.} Le \omterm{def:b1:electronic-effects:intermediates}{carbocation} est plan et \omterm{def:b1:stereochemistry-in-depth:stereocentre}{achiral}~; l'eau s'additionne
également sur les deux faces~: alcool racémique.
\textbf{18.} Une première barrière élevée (ionisation) vers le
\omterm{def:b1:electronic-effects:intermediates}{carbocation}, situé dans un puits peu profond, puis une barrière basse
vers la capture~: la première étape est cinétiquement déterminante.
\textbf{19.} Il formerait un \omterm{def:b1:electronic-effects:intermediates}{carbocation} primaire (trop instable), et
l'eau est un \omterm{def:b1:electronic-effects:nucleophile}{nucléophile} trop faible pour une SN2 rapide.
\textbf{20.} $k = A\,e^{-E_a/RT}$.
\textbf{21.} $E_a = R\ln(k_2/k_1)/(1/T_1 - 1/T_2)$.
\textbf{22.} $E_a = 8.314 \times \ln(9.74/3.00)/(1/298.15 - 1/308.15) =
\qty{9.0e4}{J/mol}$.
\textbf{23.} La hauteur de la barrière de l'ionisation, \omterm{def:b1:elementary-steps:rds}{étape
cinétiquement déterminante}.
\textbf{24.} $k = 3.0 \times 10^{-4} \exp\bigl(\frac{90\,000}{8.314}
(\frac{1}{298.15} - \frac{1}{318.15})\bigr) = \qty{2.9e-3}{s^{-1}}$.
\textbf{25.} \textbf{$E_a = \qty{90}{kJ/mol}$}.
\end{solution}
